\chapter{Appendix 13: Scenario Nu - Circular Textiles}

\section{The Legacy Paradigm \& Its Failures}
Fast fashion is driven by ecologically devastating extraction (cotton monocultures, petrochemical synthetics) and planned obsolescence. Garments are treated as disposable commodities, and the legacy consumer lacks the skills or tools to repair them, resulting in massive landfill waste and watershed microplastic pollution.

\section{The Incremental Automation Pathway}
\textbf{Phase 1: Manual Trust.} Citizens seek out local tailors through ad-hoc networks, manually negotiating fiat or barter for repairs.

\textbf{Phase 2: AI-Assisted Policy.} A BPMN engine matches repair bounties with certified tailors and tracks the lifespan of garments.

\textbf{Phase 3: Autonomous Cryptographic Enforcement.} Layer 5 mathematically incentivizes repair over replacement via escrowed Value Tokens, rigorously routes synthetic scraps to 3D printing foundries, and algorithmically restricts physical power to industrial sewing machines based on Verifiable Credentials.

\section{The Human Boundary (Limits of Automation)}
The aesthetic and structural quality of visible mending (e.g., Sashiko repair) is deeply human. While the system can restrict machine access and route scraps, the tactile feedback of fabric and the artistry of the mend are inherently manual and subjective.

\section{Orchestration Mechanics (The "How")}
The system evaluates a \texttt{RepairBounty} against the physical inventory and the DID credentials of the tailor. L5 policy gates (Chapter 4) interact with L2 smart plugs to physically restrict machine operation. Upon physical handoff (NFC logged), the multi-sig escrow (Chapter 4) automatically settles the Value Token transaction.
